\section{Conclusion}
L’analyse des travaux antérieurs montre que la gestion durable de l’arrosage nécessite la combinaison des données météorologiques, de l’humidité  du sol et des caractéristiques des cultures. Les méthodes basées sur l'évapotranspiration, notamment FAO-56 Penman-Monteith \cite{allen_crop_1998-1} et le coefficient Kc, constituent des bases pour estimer les besoins en eau. Toutefois, leur adaptation aux conditions locales, ainsi que la fiabilité et le coût des capteurs, restent des défis importants. La salinité des sols constitue également une contrainte majeure nécessitant des stratégies de correction adaptées. Ces constats justifient ainsi la mise en place d’une solution numérique intégrée, adaptée au contexte local, pour améliorer la gestion de l’irrigation et favoriser une agriculture plus durable.